\documentclass[11pt,a4paper]{article}
\usepackage[margin=24mm]{geometry}
\usepackage[T1]{fontenc}
\usepackage{lmodern,amsmath,amssymb,booktabs,tabularx,hyperref}
\hypersetup{colorlinks=true,urlcolor=blue,citecolor=blue}
\title{Device-Aware Low-Order Floating-Point Coupling\\for STT-MRAM p-Bit Ising Machines}
\author{Du Yanzhang\\National University of Singapore\\Supervisor: Prof. Kelvin Fong}
\date{Weighted Max-Cut benchmark migration --- 6 November 2026}
\begin{document}
\maketitle
\begin{abstract}
We study low-order floating-point (FP) Ising coupling representations
under nonideal stochastic-MTJ p-bit responses, compared with carefully
calibrated integer (INT) and shared-exponent block-floating (BFP) controls.
The active application benchmark is weighted Max-Cut on unmodified
Mac instances from the Biq Mac Library, as selected with the supervisor.
Graph weights are not altered to create an artificial FP advantage.
This is a research protocol, not evidence of an implemented Max-Cut
solver or physically measured p-bit. Prior MIMO work is archived.
\end{abstract}
\section{Research question}
Which coupling representation and scale grouping preserves useful
original-weight Max-Cut solution quality under a documented nonideal
STT-MRAM p-bit response with lower full-architecture hardware cost?
Novelty requires a tested configuration rule and implementable
architecture; arithmetic quantization alone is not a novelty claim.
\section{Exact objective and Ising convention}
For an undirected graph with original signed edge weights $w_{ij}$,
one vertex is one spin $s_i\in\{-1,+1\}$. Define
\begin{align}
C_w(\mathbf{s})&=\sum_{i<j}w_{ij}\frac{1-s_is_j}{2},\\
J_{ij}&=-w_{ij},\quad J_{ji}=J_{ij},\quad J_{ii}=0,\quad h_i=0,\\
E(\mathbf{s})&=-\sum_{i<j}J_{ij}s_is_j
               =W-2C_w(\mathbf{s}),\quad W=\sum_{i<j}w_{ij},\\
H_i&=\sum_{j\ne i}J_{ij}s_j,\qquad \Delta E_i=2s_iH_i .
\end{align}
This convention makes minimizing $E$ equivalent to maximizing the
original cut. Native Max-Cut needs no auxiliary spins.
Mathematical $h_i=0$ does not mean that a physical p-bit has zero
offset. Both original and quantized Hamiltonians must preserve symmetry.
\section{Source data and experimental units}
The official Biq Mac Library separates Biq quadratic programs from
Mac Max-Cut graphs; use only the Mac part
(\url{https://biqmac.aau.at/biqmaclib.html}).
Its Rudy graph format specifies an initial $n,m$ and exactly $m$
one-based weighted edges. The initial source-audit shortlist is
$w01,w05,w09$ (signed), $pw01,pw05,pw09$ (nonnegative), and
$pm1s,pm1d$ (signed sparse/dense), with $n=100$ and instance
suffix .0. Preserve raw files and SHA256. Later experiments may
split seeds .0--.4 for calibration and .5--.9 for held-out evaluation;
results on a single .0 graph cannot establish a trend.
\section{Measurements and evidence}
Score states against the \emph{original source weights}, not the
represented Hamiltonian: original cut value, absolute difference to
a genuinely certified optimum, or separately identified bounds.
A returned state must come from a declared implementable policy
(such as lowest visited \emph{represented} energy), not an uncosted
oracle using the original objective. Use paired graph, solver seed and
device population with a fixed attempted-update budget.
Visited-target probability is distinct from returned-target success.
Do not call an upper-bound difference the exact optimality gap.
Negative edge weights and zero optima make some relative ratios
inappropriate.

Evidence categories are E0 exact mapping, E1 seeded algorithmic
performance, E2 illustrative response, E3 sourced/validated stochastic
device transfer, E4 authorized circuit simulation, and E5 RTL/synthesis.
These cannot be interchanged.
\section{Representation and physical-device design}
The project retains low-order FP4/FP6/FP8 couplings, with optimized
scaled INT and BFP as fair controls. Many Mac/Rudy source weights are
small integers: under unit scale, signed six-bit INT/BFP can represent
the frequently used $[-10,10]$ weights exactly, whereas FP4 need not.
Therefore FP superiority must be established experimentally or rejected.
Fair comparisons include coefficient bits, global/group metadata,
decoder, scale alignment, arithmetic width, memory reads, state
controller, device interface and timing constraints.
The original problem bias is zero, so bias storage and circuitry must
be removed or tied off \emph{before} claiming Max-Cut-specific area.

The digital local-field row is not automatically an in-array
MRAM-CIM macro. The magnetic p-bit response, including fixed offsets,
gain, correlation, voltage and timescale, requires physically
grounded device evidence. Existing stable-storage GF22FDX MTJ models
are not proof of a stochastic low-barrier p-bit.
\section{Engineering evidence retained from the prior workload}
A cycle-serial local-field row previously passed 84/84 Icarus
bit-true vectors, and three bias-inclusive Nangate45/Yosys
mappings; the standalone best-represented-energy tracker passed
a focused RTL test and mapped to 3265.150 library-area units.
These engineering tests are retained as \emph{component} evidence,
not reclassified as weighted-Max-Cut quality or full-system area.
A previous wired row/tracker wrapper has not yet been
reported as successfully NUS-tested.
\section{Stage gates}
\begin{center}\small
\begin{tabularx}{\linewidth}{@{}lX@{}}\toprule
M0 & Import original Biq Mac Mac graphs; SHA256 and E0
cut--Ising--flip checks; no fabricated optimum.\\
M1 & Original-cut scoring and frozen, paired FP/INT/BFP baseline
with honest returned-solution policy.\\
M2 & Zero-bias device-aware coupling design rule and cost model.\\
M3 & Bit-true Max-Cut row/tracker RTL integration and mapped
area/timing with memory/peripheral accounting.\\
M4 & Held-out signed-weight/density/size and device-model tests.\\
M5 & Qualified final design rules, limitations, and report.\\\bottomrule
\end{tabularx}
\end{center}
\section{Current status}
The new Rudy parser, mathematical checker and tests were
authored in the M0 draft branch. They are not yet NUS-executed.
No new Biq Mac optimization quality, measured stochastic-MTJ
device response, or Max-Cut-specific mapped-area claim exists.
The previous MIMO-oriented report is preserved at
\texttt{docs/archive/MAIN\_MIMO\_PRE\_MAXCUT\_20261106.tex}.
\end{document}
